\documentclass{article}
\usepackage{graphicx} % Required for inserting images

\title{aaaaa}
\author{Jordi Rodrígo Lozada Rivera}
\date{August 2026}

\begin{document}

\begin{abstract}
En esta práctica se determinó experimentalmente la relación entre la constante de
Planck y la carga elemental, $h/e$, utilizando diodos emisores de luz (LED) de
diferentes longitudes de onda. Para cada LED se midió el voltaje ánodo--cátodo
correspondiente al umbral de emisión luminosa y, a partir de la longitud de onda
nominal, se calculó la frecuencia de la radiación emitida. Se construyó la gráfica
de frecuencia en función del voltaje y se realizó un ajuste lineal cuya pendiente
permitió estimar la relación $h/e$. El valor experimental obtenido fue: 

\[
\left(\frac{h}{e}\right)_{\mathrm{exp}}
=(3.992\pm0.059)\times10^{-15}\ \mathrm{V\,s},
\]

con un error relativo de $3.47\,\%$ respecto al valor aceptado. Los resultados
confirman la relación lineal predicha por el modelo del efecto fotoeléctrico y
muestran que el empleo de LEDs constituye un método sencillo y confiable para
estimar constantes fundamentales.
\end{abstract}


\textbf{Palabras clave:}
Efecto fotoeléctrico, LED, constante de Planck, carga del electron.

\section{Introducción}

El efecto fotoeléctrico constituye uno de los experimentos fundamentales de la
física moderna, pues demuestra que la energía de la luz está cuantizada en
fotones de energía $E=h\nu$, donde $h$ es la constante de Planck y $\nu$ la
frecuencia de la radiación. La explicación propuesta por Albert Einstein en
1905 permitió interpretar diversos fenómenos que no podían describirse mediante
la teoría clásica del electromagnetismo y representó uno de los primeros
evidencias de la naturaleza cuántica de la luz. \\

Un diodo emisor de luz (LED) funciona mediante la recombinación de electrones y
huecos en un material semiconductor. Durante este proceso se emiten fotones cuya
energía está relacionada con la diferencia de energía entre las bandas del
material. Como primera aproximación puede considerarse que la energía adquirida
por un electrón al atravesar la unión del diodo es igual a la energía del fotón
emitido,

\[
eV=h\nu,
\]

donde $e$ es la carga elemental y $V$ es el voltaje umbral del LED. Esta
expresión predice una relación lineal entre la frecuencia de la luz emitida y el
voltaje aplicado, cuya pendiente corresponde a $e/h$.\\

En esta práctica se utilizaron LEDs de diferentes longitudes de onda para medir
su voltaje de umbral y construir la gráfica de frecuencia contra voltaje.
Mediante un ajuste lineal se obtuvo una estimación experimental de la relación
$h/e$, comparándola posteriormente con el valor aceptado y evaluando el error
porcentual del método.

\section{Discusión y Conclusión}

Los resultados obtenidos muestran una relación aproximadamente lineal entre la
frecuencia de emisión de los LEDs y el voltaje ánodo--cátodo medido en el umbral
de conducción. El ajuste lineal presentó un coeficiente de determinación
$R^2=0.99089$, indicando que el modelo empleado describe adecuadamente el
comportamiento experimental.\\

A partir de la pendiente del ajuste se obtuvo

\[
\left(\frac{h}{e}\right)_{\mathrm{exp}}
=(3.992\pm0.059)\times10^{-15}\ \mathrm{V\,s},
\]

valor que difiere únicamente un $3.47\,\%$ del valor aceptado
$4.135668\times10^{-15}\ \mathrm{V\,s}$. Esta diferencia resulta razonable para
un experimento de laboratorio y confirma que el procedimiento permite estimar
adecuadamente la relación entre ambas constantes fundamentales.\\

Las principales fuentes de incertidumbre provienen de la determinación del
voltaje umbral de cada LED, la dispersión de las longitudes de onda nominales
proporcionadas por el fabricante, la resolución de los instrumentos de medición
y la aproximación de considerar que toda la energía eléctrica suministrada se
transforma en energía luminosa. En realidad, parte de la energía se disipa como
calor dentro del semiconductor, lo cual introduce una desviación sistemática
respecto al modelo ideal.\\

En conclusión, la práctica permitió verificar experimentalmente la dependencia
lineal entre la frecuencia de la luz emitida y el voltaje umbral de los LEDs,
obteniendo una estimación de la relación $h/e$ con un error menor al $4\,\%$.
Los resultados respaldan el modelo cuántico de emisión de luz y muestran que el
uso de LEDs constituye una herramienta didáctica sencilla para determinar
constantes fundamentales mediante técnicas experimentales accesibles.

\section{Conclusiones}

La práctica permitió verificar experimentalmente la dependencia
lineal entre la frecuencia de la luz emitida y el voltaje umbral de los LEDs,
obteniendo una estimación de la relación $h/e$ con un error menor al $4\,\%$.
Los resultados respaldan el modelo cuántico de emisión de luz y muestran que el
uso de LEDs constituye una herramienta didáctica sencilla para determinar
constantes fundamentales mediante técnicas experimentales accesibles.
\end{document}
